\section{Prospective selection with generated HTTP responses}
\label{sec:realapi}

\subsection{Protocol and candidate pool}
This study tests the same prediction--selection--execution contract with actual
response generation at the offload endpoint. On one host with four A800 80GB
PCIe GPUs, two TP1 Qwen2.5-7B-Instruct FP16 replicas are elastic local workers;
two further replicas remain ready as HTTP providers. This is a two-local,
two-provider topology, distinct from Section~Q's four elastic workers. The
stack is vLLM 0.9.0, PyTorch 2.7.0+cu126, and Transformers 4.51.1.
All replicas use the same model weights and tokenizer. The recovery and steady
inputs are the previously locked deployment windows: 358/944 online requests
and 120 offline jobs each, with unchanged token IDs, arrivals, and deadlines.
These are known reused inputs, not newly held-out arrivals.

The candidate pool contains H, H+guard, PPO-1--5, and the earlier mixed-policy
PPO-C (displayed here as PPO-12). All six checkpoints already exist; none is
retrained, selected from new GPU outcomes, or expanded to four-device actions.
The matrix is eight candidates times two windows times two repeats, plus
four all2 references: 36 executions. All2 keeps both local replicas ready;
the same two provider replicas remain allocated for every cell. Each run has
1,800~s of arrivals and a 300~s drain. The strict P1+ test includes every
planned request, including late, undispatched, and unfinished work.

Frozen preparation binds inputs, checkpoint and source hashes, profiles,
72 predictions, and selected sets before the first formal setup. The 72
predictions comprise four predictors times eight candidates times two windows,
plus eight reference predictions. The first formal setup begins 16.011~s
after the freeze; all 36 execution receipts reference the same preparation
and source seals. There are no formal retries or excluded scientific failures.
The two physical repeats reuse each prediction, rather than constituting
independent predictor samples.

\subsection{HTTP execution and calibration}
Local-ready slots retain online-first dispatch unless the guard reserves an offline slot. Online requests still queued
10~s after arrival may offload to eight fixed provider slots (four per
provider); offline jobs cannot offload. Each call submits the registered token
IDs directly, without chat-template insertion or retokenization. Performance
runs use greedy decoding with seed zero, ignored EOS, and the registered
128/256 online/offline output budgets. HTTP completion requires a successful,
validated generated response. There is no five-second padding or application
retry. Native provider responses and client receipts are matched by request
ID, payload hash, response hash, usage, and monotonic timestamps.

Independent profiling precedes the freeze: five input lengths
(128, 256, 512, 1024, 2048), two output budgets (128, 256), three concurrency
levels (1, 2, 4), three batches, and four GPUs give 360 batches. Replay uses
the median batch maximum, interpolated in input length and concurrency.
Remaining service work advances at the rate for the current active
concurrency. HTTP proxy overhead is a separate terminal stage. Formal response
latencies do not fit these profiles.

Twenty launch-to-ready samples, with another local replica continuously
serving four background jobs, give a median of 29.195~s, p95 of 29.384~s,
and maximum of 29.447~s. Local uses the measured median. The 69~s and 30~s
predictors replace only startup with 69.219~s or 30~s, holding service,
shutdown, and guard estimates fixed. These names denote study-specific
models, not the original A100 parameter table: the 69~s variant carries over
only the archived A100 startup estimate without remeasurement on A800.
A fourth, CPU-only diagnostic
retains Local's lifecycle but replaces provider service with an eight-slot
five-second timer. All predictors and physical runs use the same separately
frozen guard safety profile and five-second margin.

\subsection{Accounting and selection}
The primary system ledger charges local occupation and both fixed providers:
\begin{equation}
 C_{\rm sys}=.001(G_{\rm local}+4200)+.002S+.001D.
\end{equation}
The 4,200 GPU seconds are two providers allocated for the full 2,100~s
observation, including when idle. Local GPU occupation extends from launch to
verified release, clipped to the observation interval. Setup and cleanup are
separate. Feasibility gates selection, so a deadline penalty is unnecessary
in this cost. Values are scenario USD, not measured cloud bills.

A secondary client ledger charges local occupation and lifecycle events plus
$10^{-6}I+3\times10^{-6}B$ for token usage in successful HTTP responses.
It describes a different pricing scenario; it is neither added to the system
ledger nor used for the primary selection. In particular, provider capacity
is not treated as free in the primary result.

For each predictor, eligibility requires predicted P1+ in both windows.
Selection minimizes equal-window mean predicted system cost, retaining all
ties within $10^{-9}$. Physical comparison requires P1+ in all four
candidate executions, then averages the two windows and two repeats equally.
All2 is a capacity reference outside the eight-candidate selection pool.

\begin{table*}[t]
\centering\small
\caption{All prospective physical outcomes. R/S list recovery/steady P1+
passes out of two repeats. Costs are system scenario USD, repeat 1 / repeat 2;
$\dagger$ marks an infeasible execution. Every online request meets its deadline.}
\label{tab:realapi-all}
\setlength{\tabcolsep}{4pt}
\begin{tabular*}{\textwidth}{@{\extracolsep{\fill}}lccrrr@{}}
\toprule
Candidate & R & S & Recovery costs & Steady costs & Mean cost\\
\midrule
H & 1/2 & 1/2 & 5.1637 $\dagger$/5.0278 & 5.3513/5.3235 $\dagger$ & 5.2166\\
H+guard & 2/2 & 2/2 & 5.0281/5.2124 & 5.2251/5.3644 & 5.2075\\
PPO-1 & 2/2 & 2/2 & 6.2337/6.2326 & 6.2507/6.2514 & 6.2421\\
PPO-2 & 2/2 & 2/2 & 5.3627/5.3646 & 5.3680/5.3691 & 5.3661\\
PPO-3 & 2/2 & 2/2 & 5.3657/5.3706 & 5.3706/5.3622 & 5.3673\\
PPO-4 & 2/2 & 2/2 & 6.1432/5.4045 & 5.3805/5.3744 & 5.5756\\
PPO-5 & 2/2 & 2/2 & 5.9622/5.8927 & 5.8892/5.8902 & 5.9086\\
PPO-12 & 2/2 & 2/2 & 5.8054/5.8041 & 5.6964/5.7007 & 5.7516\\
all2 & 2/2 & 2/2 & 8.4000/8.4000 & 8.4000/8.4000 & 8.4000\\
\bottomrule

\end{tabular*}
\end{table*}

\subsection{Physical results and choice consequences}
All 36 cells are technically valid; 34 satisfy strict P1+. Both negative
cells are H: recovery repeat 1 and steady repeat 2 each leave the last
offline job unfinished. Its arrival is 1,785~s and deadline 2,085~s. It remains
undispatched through the cutoff, with one offline request queued, no online
queue, both local workers off, and H's target zero. These are local lifecycle
and admission failures, not HTTP failures. H passes its other two executions.
H+guard and every PPO checkpoint pass all four executions, giving seven
consistently feasible candidates (Table~\ref{tab:realapi-all}).

Across all cells, 11,137 offloads have successful validated responses:
5,577 and 5,560 on the two providers. There are no HTTP errors, missing usage,
or responses completed after the observation cutoff. Total recorded provider
usage is 2,286,615 input tokens and 1,425,536 output tokens. Both endpoints
serve offloads for every candidate. These responses establish actual inference
on the offload path, without implying independent network or model diversity.

Local selects H+guard, whose mean system cost is 5.2075. It is the cheapest
of the seven consistently feasible candidates both overall and in each
repeat's equal-window comparison. The historical startup proxy rejects the
guard on steady and selects PPO-1, costing 6.2421. Its relative premium is
$(6.2420812-5.2075119)/5.2075119=19.87\%$; the equivalent reduction relative
to PPO-1 is 16.57\%. This is a within-pool, fixed-provider-allocation consequence,
not a new policy-optimization result or a commercial API saving.

\begin{table}[t]
\centering\small
\caption{Prediction versus 32 candidate executions. FA is out of two
infeasible cells; FR out of 30 feasible cells. Each predictor has 16 distinct
candidate/window predictions, reused across repeats. Cost MAE is system scenario USD.}
\label{tab:realapi-predictions}
\setlength{\tabcolsep}{3pt}
\begin{tabular*}{\linewidth}{@{\extracolsep{\fill}}lrrrr@{}}
\toprule
Predictor & Agreement & FA & FR & Cost MAE\\
\midrule
Local & 30/32 & 1/2 & 1/30 & 0.0610\\
69\,s & 16/32 & 1/2 & 15/30 & 0.4661\\
30\,s & 30/32 & 1/2 & 1/30 & 0.0610\\
Local + timer & 30/32 & 1/2 & 1/30 & 0.0644\\
\bottomrule
\end{tabular*}
\end{table}

\begin{table}[t]
\centering\small
\caption{Key frozen predictions behind the prospective choice. R/S count
on-time offline jobs out of 120; every predicted online request is timely.
Cost is the predicted two-window mean. An ineligible candidate is excluded
before cost ranking.}
\label{tab:realapi-eligibility}
\setlength{\tabcolsep}{3pt}
\begin{tabular*}{\linewidth}{@{\extracolsep{\fill}}llrrrl@{}}
\toprule
Replay & Candidate & R & S & Cost & Eligible\\
\midrule
Local & H+guard & 120 & 120 & 5.1394 & yes\\
Local & PPO-1 & 120 & 120 & 6.2290 & yes\\
69\,s & H+guard & 120 & 117 & 6.0235 & no\\
69\,s & PPO-1 & 120 & 120 & 6.1757 & yes\\
\bottomrule
\end{tabular*}
\end{table}

The 69~s model predicts three late offline completions on steady, excluding
H+guard despite predicting a lower mean cost than PPO-1
(Table~\ref{tab:realapi-eligibility}). The physical guard has no misses. This
choice difference follows the feasibility gate, not a cost tie. The traces
support that selection consequence; we do not attribute this entire HTTP
trajectory to one isolated startup event.

Local is not a perfect classifier (Table~\ref{tab:realapi-predictions}). It
accepts H on recovery, falsely accepting repeat 1, and rejects H on steady,
falsely rejecting repeat 1. Rejecting H on steady nevertheless removes it
from two-window selection. The 30~s predictor has identical predictions here;
the timer diagnostic also selects H+guard and has the same labels. Therefore,
this study validates a frozen selection with generated responses but does
not show that local calibration or HTTP service modeling is necessary to
select the winner on these inputs. It complements the earlier failure cases
with a prospective successful choice and a measured alternative-choice cost.

\subsection{Paired endpoint output check}
A separate preformal check uses 64 hash-selected GSM8K test items
\citep{cobbe2021gsm8k}, submitted in pairs through temporary local-HTTP and
provider-HTTP replicas with the same weights, prompt template, and greedy
decoding. This probe permits natural EOS with a maximum of 512 output tokens;
it is separate from the fixed-budget performance inputs and does not select
policies. Both routes answer 61/64 correctly, with the same per-item correctness
vector, three shared failures, including one unparsed answer per route.
There are no truncated responses or HTTP errors. The local-HTTP replica is
not the formal local worker implementation. These descriptive observations
check the configured endpoints; they neither establish statistical quality
equivalence nor characterize heterogeneous commercial models.

\subsection{Portable evidence and audit scope}
The anonymous repository's \path{paper/artifact/realapi_prospective/verify.py} extracts the sealed
anonymous evidence into a temporary directory and runs its standard-library
CPU verifier. It checks all 36 normalized physical records, actual HTTP
response hashes and usage, 64 candidate predictions, frozen selected sets,
and all 128 paired quality responses. It also recomputes table means,
feasibility counts, classification errors, and selected-policy cost gaps.
The eight frozen all2 predictions are outside the released candidate-prediction
set; all four physical references are included. These checks audit recorded
execution and arithmetic; they do not rerun model inference or authenticate
hardware independently. Full native journals and calibration provenance are
retained separately from this portable release.
